%% kapa: exam=1 round=19 year=2008 period=1 form=A subject=civil_law subject_ko=민법 seq=15 number=15
%% source: Q-Net 감정평가사 기출문제 (cst00309) artlSeq=1001962 file=extracted/1008/1001962/2008년 제19회 감정평가사 1차시험 기출문제/1교시.hwp
%% conversion: kordoc 4.24.51 -> tools/gen1.py (auto); review=proofread-sonnet; flags=none
\begin{question}{15}
다음 사례에 관한 설명으로 틀린 것은? (다툼이 있는 경우에는 판례에 의함)
\begin{framedtext}
A. 甲은 乙과 도박을 하던 중 돈을 모두 잃었다. 그러자 甲은 乙에게 100만원을 도박자금으로 빌려달라고 하였고, 이에 따라 乙은 甲에게 100만원을 빌려주었다.\par
B. 丙은 丁에 대한 도박채무를 변제하기 위하여 丁에게 자신의 부동산을 처분하여 그 대금을 변제에 충당하라고 하였다. 이에 丁은 丙을 대리하여 그 부동산을 선의의 제3자에게 매도하고 그 대금으로 변제에 충당하였다.
\end{framedtext}
\begin{choices}
  \item A의 경우, 甲과 乙사이의 금전소비대차계약은 무효이다.
  \item A의 경우, 乙은 甲에 대하여 대여금의 반환을 청구할 수 없다.
  \item B의 경우, 丙의 丁에 대한 도박채무의 부담행위 및 그 변제약정은 선량한 풍속 기타 사회질서에 위반된다.
  \item B의 경우, 제3자는 부동산에 대한 소유권을 취득할 수 있다.
  \item B의 경우, 도박채무를 이행하기 위해 丁에게 대리권을 수여한 행위는 무효이다.
\end{choices}
\end{question}
\answer{5}
